\documentclass[11pt,a4paper]{article}
\usepackage[T1]{fontenc}
\usepackage[utf8]{inputenc}
\usepackage[main=italian,provide=*]{babel}
\usepackage{amsmath,amssymb}
\usepackage{geometry}
\geometry{margin=2.3cm,top=2.5cm,bottom=2.7cm}
\usepackage{booktabs}
\usepackage[table,dvipsnames]{xcolor}
\usepackage{tcolorbox}
\tcbuselibrary{skins,breakable}
\usepackage{titlesec}
\usepackage{enumitem}
\usepackage{needspace}
\usepackage{fancyhdr}
\usepackage{graphicx}
\usepackage{hyperref}
\hypersetup{colorlinks=true,linkcolor=Sepia,citecolor=Sepia,urlcolor=Sepia}
\newcommand{\Tr}{\operatorname{Tr}}
\newcommand{\ket}[1]{\lvert #1 \rangle}
\newcommand{\bra}[1]{\langle #1 \rvert}

\definecolor{inkblue}{HTML}{1B3A5C}
\definecolor{lightblue}{HTML}{EAF1F8}
\definecolor{accent}{HTML}{B0562A}
\definecolor{softgray}{HTML}{6B6B6B}
\definecolor{okgreen}{HTML}{2E6B3E}
\definecolor{okgreenbg}{HTML}{EAF5EC}

\titleformat{\section}
  {\normalfont\Large\bfseries\color{inkblue}}
  {\thesection}{0.6em}{}[{\color{inkblue!40}\titlerule}]
\titlespacing{\section}{0pt}{0.9em}{0.4em}
\titleformat{\subsection}
  {\normalfont\large\bfseries\color{accent}}
  {}{0em}{}
\titlespacing{\subsection}{0pt}{0.6em}{0.2em}

\pagestyle{fancy}
\fancyhf{}
\renewcommand{\headrulewidth}{0.4pt}
\fancyhead[L]{\small\color{softgray}Il modello di rumore per il dimero}
\fancyhead[R]{\small\color{softgray}\thepage}
\fancyfoot[C]{\small\color{softgray}Tesi triennale, Samuele --- Relatore: Prof.\ Alessandro Chiesa}

\newtcolorbox{keyeq}[1][]{
  colback=lightblue, colframe=inkblue!70, boxrule=0.6pt,
  arc=2pt, left=8pt, right=8pt, top=3pt, bottom=3pt, #1
}
\newtcolorbox{okbox}[1]{
  colback=okgreenbg, colframe=okgreen!60, boxrule=0.6pt,
  arc=2pt, left=7pt, right=7pt, top=2pt, bottom=2pt,
  fonttitle=\bfseries\color{okgreen}, title={#1}
}

\title{\vspace{-1.5em}\color{inkblue}Il modello di rumore per il dimero\\[0.2em]
\Large\normalfont\color{softgray}Canale depolarizzante, errore di lettura, calibrazione reale}
\author{Note di lavoro --- Tesi triennale, Samuele\\ Relatore: Prof.\ Alessandro Chiesa}
\date{}

\begin{document}
\sloppy
\maketitle
\thispagestyle{fancy}
\vspace{-2.0em}

File di codice: \texttt{noise\_model\_dimero.py}. Derivazione completa in
\texttt{teoria\_modello\_rumore.tex} (qui richiamate solo le formule
essenziali). Nessun punto di lavoro specifico: questo è il modello di
rumore condiviso da tutti gli stadi successivi della pipeline.

\tableofcontents

\section{Obiettivo}

Prima di poter dire quanto un risultato VQE, una simulazione di Trotter o
un correlatore dinamico degradino sotto rumore, serve un modello di
rumore concreto: quali canali includere, con quali parametri, calibrati
su cosa. Questo documento fissa quella base, condivisa da tutti gli
stadi a valle (VQE con termine DM sotto rumore, quantum simulation
Trotter sotto rumore, correlatori dinamici sotto rumore, scan sui
parametri di rumore).

Il relatore ha confermato esplicitamente (\texttt{domande\_relatore.md},
sez.~12) lo scope: \textbf{due soli canali}, errore di gate
(depolarizzante, 1 e 2 qubit) ed errore di lettura (readout, simmetrico
come primo passo). $T_1/T_2$ esclusi per decisione esplicita --- non
un'omissione, una scelta di scope dichiarata.

\section{Il canale depolarizzante: derivazione essenziale}

Un canale depolarizzante a $n$ qubit ($d=2^n$) sostituisce lo stato con
una miscela fra lo stato originale e lo stato massimamente misto:
\begin{keyeq}
\centering
$\mathcal{E}(\rho) = (1-\lambda)\,\rho + \lambda\,\dfrac{\mathbf{1}}{d}$
\end{keyeq}
dove $\lambda\in[0,1]$ è il parametro del canale in Qiskit Aer. Le
calibrazioni pubblicate (randomized benchmarking) non riportano $\lambda$
direttamente, ma l'\textbf{errore medio di gate} $\varepsilon$ ---
la probabilità media che il gate produca un risultato diverso da quello
ideale, mediata su tutti gli stati di ingresso. Il legame fra i due,
derivato in \texttt{teoria\_modello\_rumore.tex}, è
\begin{keyeq}
\centering
$\varepsilon = \lambda\,\dfrac{d-1}{d} \quad\Longrightarrow\quad
\lambda_{1q} = 2\,\varepsilon_{1q} \ \ (d=2), \qquad
\lambda_{2q} = \dfrac{4}{3}\,\varepsilon_{2q} \ \ (d=4)$
\end{keyeq}

\begin{figure}[h]
\centering
\includegraphics[width=0.62\linewidth]{figure/fig_depolarizzazione_bloch.pdf}
\caption{\textit{Contrazione del vettore di Bloch, $1-\lambda$, in
funzione di $\lambda$. Ai valori di calibrazione di riferimento (linee
tratteggiate) la contrazione è piccola ($<1\%$) ma non nulla --- l'intero
punto della pipeline di rumore è misurare quanto questo piccolo effetto,
accumulato su molti gate, degradi i risultati a valle.}}
\end{figure}

\section{Errore di lettura}

Il readout simmetrico è descritto da un solo parametro $p$: la
probabilità di leggere il risultato sbagliato è la stessa in entrambe le
direzioni, $0\to1$ e $1\to0$. La correzione sull'osservabile
$\langle Z\rangle$ --- derivata per sostituzione diretta nella
distribuzione di probabilità letta, dettagli in
\texttt{teoria\_modello\_rumore.tex} --- è
\begin{keyeq}
\centering
$\langle Z\rangle_\text{readout} = (1-2p)\,\langle Z\rangle_\text{ideale}$
\end{keyeq}
un fattore puramente moltiplicativo. (Il caso $p_{01}\neq p_{10}$,
asimmetrico, è trattato separatamente in
\texttt{teoria\_readout\_asimmetrico.tex}: introduce un termine additivo
che complica questa semplice proporzionalità.)

\section{Calibrazione: valori di riferimento}

\begin{keyeq}[title={Fonte citata}]
\centering
J.~P.~T.~Stenger, G.~Bazargan, N.~T.~Bronn, D.~Gunlycke, \emph{``Method
for simulating open-system dynamics using mid-circuit measurements on a
quantum computer''}, arXiv:2504.15187 (2025) --- Appendice B, mediane
pubblicate per \texttt{ibm\_torino} (IBM Heron r1, 133 qubit).
\end{keyeq}

\begin{center}
\renewcommand{\arraystretch}{1.3}
\begin{tabular}{@{}llc@{}}
\toprule
\textbf{Canale} & \textbf{Gate di riferimento} & \textbf{Valore mediano}\\
\midrule
Errore di gate, 1 qubit & \texttt{sx} & $\varepsilon_{1q}=2.9\times10^{-4}$\\
Errore di gate, 2 qubit & \texttt{CZ} & $\varepsilon_{2q}=3.8\times10^{-3}$\\
Errore di lettura & --- & $p_\text{readout}=2.3\times10^{-2}$\\
\bottomrule
\end{tabular}
\end{center}

\begin{okbox}{Nota di modellazione, dichiarata esplicitamente}
Il gate nativo a 2 qubit di \texttt{ibm\_torino} è il CZ, mentre il
circuito di questa tesi è scritto nella base $\{rz,sx,x,cx\}$. Si usa il
valore di errore del CZ come rappresentativo dell'ordine di grandezza
dell'errore di gate a 2 qubit --- lo stesso approccio generico richiesto
dal relatore (\emph{"usa dei valori tipici... guardando qualche chip di
riferimento"}), non una replica esatta dell'hardware. Lo scan sui
parametri di rumore copre comunque un intorno di questi valori, non solo
il punto singolo.
\end{okbox}

\subsection{Perché il valore del CZ è una scelta giustificata, non solo tipica}

L'argomento è più solido di un semplice ``stesso ordine di grandezza fra
gate simili'': il CNOT e il CZ sono legati da un'identità circuitale
\textbf{esatta}, non da un'approssimazione,
\begin{keyeq}
\centering
$\text{CNOT} = (\mathbb{1}\otimes H)\cdot\text{CZ}\cdot(\mathbb{1}\otimes H)$
\end{keyeq}
Il CZ è il gate nativo di \texttt{ibm\_torino}: fisicamente, l'hardware
realizza un'interazione a due qubit che \emph{è} il CZ, non il CNOT.
Quando il transpilatore deve produrre un CNOT logico su questo hardware,
genera esattamente questa sequenza --- due Hadamard (gate a 1 qubit,
gia' inclusi separatamente in $\varepsilon_{1q}$) attorno al CZ nativo.
La parte difficile, l'interazione a due qubit dove si concentra la
maggior parte dell'errore, è la \emph{stessa identica operazione fisica}
in entrambi i casi --- non un'analogia fra gate della stessa famiglia,
lo stesso gate con un cambio di base attorno.

\paragraph{Il calcolo, per esteso.} Nella base
$\{\ket{00},\ket{01},\ket{10},\ket{11}\}$,
\begin{keyeq}
\centering
$\text{CZ} = \begin{pmatrix}1&0&0&0\\0&1&0&0\\0&0&1&0\\0&0&0&-1\end{pmatrix},
\qquad
\mathbb{1}\otimes H = \dfrac{1}{\sqrt2}\begin{pmatrix}1&1&0&0\\1&-1&0&0\\0&0&1&1\\0&0&1&-1\end{pmatrix}$
\end{keyeq}
Il prodotto $(\mathbb{1}\otimes H)\cdot\text{CZ}\cdot(\mathbb{1}\otimes H)$
agisce a blocchi $2\times2$ sul secondo qubit, condizionato al valore del
primo: sul blocco condizionato a $\ket0$ si ha $H\cdot\mathbb{1}\cdot H=\mathbb{1}$
(l'identità coniugata da $H$ resta l'identità); sul blocco condizionato a
$\ket1$ si ha $H\cdot Z\cdot H=X$ --- l'identità standard di coniugazione
di Pauli tramite Hadamard, verificabile da
$H=\frac{1}{\sqrt2}\begin{pmatrix}1&1\\1&-1\end{pmatrix}$ per sostituzione
diretta. Il risultato è
\begin{keyeq}
\centering
$(\mathbb{1}\otimes H)\cdot\text{CZ}\cdot(\mathbb{1}\otimes H)
= \begin{pmatrix}1&0&0&0\\0&1&0&0\\0&0&0&1\\0&0&1&0\end{pmatrix} = \text{CNOT}$
\end{keyeq}
--- verificato anche numericamente (moltiplicazione matriciale diretta),
scarto dalla forma esatta del CNOT: $2.2\times10^{-16}$ (precisione
macchina, zero per costruzione algebrica, non una coincidenza numerica).

\begin{okbox}{Il residuo di approssimazione, dichiarato con onestà}
Due cose che questa identità non chiude del tutto:
\begin{itemize}[leftmargin=1.3em,itemsep=0.15em]
\item il costo aggiuntivo dei due Hadamard, $\sim2\varepsilon_{1q}\approx6\times10^{-4}$,
sopra $\varepsilon_{2q}=3.8\times10^{-3}$ --- un aumento relativo
$\sim15\%$, non trascurabile ma nemmeno un ordine di grandezza;
\item il randomized benchmarking misura l'errore medio del CZ isolato:
non è un teorema che lo stesso valore si trasferisca invariato al CNOT
sintetizzato, solo un'ipotesi ragionevole data l'identità esatta sopra.
\end{itemize}

\paragraph{Da dove viene $2\varepsilon_{1q}$.} Non è una formula esatta,
è una stima al prim'ordine. Il randomized benchmarking assegna a ogni
gate una probabilità di errore $\varepsilon$ (la probabilità che il gate
applicato non sia quello ideale). Per due gate \textbf{indipendenti} in
sequenza (i due Hadamard), la probabilità che almeno uno fallisca è
\begin{keyeq}
\centering
$P(\text{almeno un errore}) = 1-(1-\varepsilon_{1q})^2 = 2\varepsilon_{1q}-\varepsilon_{1q}^2$
\end{keyeq}
Con $\varepsilon_{1q}=2.9\times10^{-4}$, il termine quadratico è
$\varepsilon_{1q}^2\approx8.4\times10^{-8}$ --- quattro ordini di
grandezza più piccolo del termine lineare $2\varepsilon_{1q}\approx5.8\times10^{-4}$,
quindi trascurabile: da qui l'approssimazione
$P(\text{almeno un errore})\approx2\varepsilon_{1q}$. Il rapporto rispetto
a $\varepsilon_{2q}=3.8\times10^{-3}$ dà l'aumento relativo dichiarato,
$5.8\times10^{-4}/3.8\times10^{-3}\approx15\%$.
\end{okbox}
La formulazione cauta (``valore rappresentativo dell'ordine di
grandezza'') resta comunque quella corretta da riportare: soddisfa il
livello di rigore richiesto dal relatore senza sovra-vendere una
giustificazione che, pur solida, non è un calcolo esatto end-to-end.

\section{Verifica}

Costruito il \texttt{NoiseModel} con \texttt{build\_noise\_model()} ai
valori di riferimento:

\begin{center}
\renewcommand{\arraystretch}{1.25}
\begin{tabular}{@{}lc@{}}
\toprule
\textbf{Quantità} & \textbf{Valore}\\
\midrule
$\varepsilon_{1q}$ & $2.9000\times10^{-4}$\\
$\varepsilon_{2q}$ & $3.8000\times10^{-3}$\\
$p_\text{readout}$ & $2.3000\times10^{-2}$\\
$\lambda_{1q}=2\varepsilon_{1q}$ & $5.8000\times10^{-4}$\\
$\lambda_{2q}=\frac{4}{3}\varepsilon_{2q}$ & $5.0667\times10^{-3}$\\
\bottomrule
\end{tabular}
\end{center}

\begin{okbox}{Verificato}
I gate di base $\{rz,sx,x\}$ portano il canale a 1 qubit ($\lambda_{1q}$);
il \texttt{cx} porta il canale a 2 qubit ($\lambda_{2q}$); la
\texttt{measure} porta l'errore di lettura --- struttura del
\texttt{NoiseModel} verificata via \texttt{noise\_model.py} stesso
(elenco delle istruzioni rumorose stampato a schermo, coincide con
quanto dichiarato: \texttt{['rz','sx','x','cx','measure']}).
\end{okbox}

\section{Cosa non entra in questo modello}

Esplicito, per non lasciarlo implicito (stesso principio applicato in
tutte le estensioni successive):
\begin{itemize}[leftmargin=1.3em,itemsep=0.2em]
\item $T_1/T_2$ (rilassamento e decoerenza continua durante il circuito)
--- decisione esplicita del relatore, non un'omissione tecnica.
\item Correlazioni fra qubit (crosstalk di gate o di lettura): ogni
canale è applicato indipendentemente qubit per qubit (o coppia per
coppia per il \texttt{cx}).
\item Asimmetria del readout ($p_{01}\neq p_{10}$): trattata come
estensione separata, non nel modello di base.
\item Deriva temporale della calibrazione: un solo punto nel tempo
(mediane pubblicate), non un modello dipendente dal momento
dell'esecuzione.
\end{itemize}

\section{Considerazioni e cosa serve al passo successivo}

Il modello è deliberatamente minimale: due canali, quattro numeri. Questa
scelta ha una conseguenza diretta sul passo successivo (VQE con termine
DM sotto rumore): l'unico modo in cui il rumore entra è tramite il
\texttt{NoiseModel} costruito qui, applicato durante la simulazione a
operatore densità --- nessun altro canale, nessun parametro nascosto.
Chi legge i risultati del passo successivo deve sapere che ``rumoroso''
significa, in questa tesi, esattamente questi due canali a questi
valori (o a quelli esplicitamente scansionati) --- non una replica
generica di ``un computer quantistico rumoroso''.

\begin{okbox}{Sintesi in una riga}
Due canali (depolarizzante 1q/2q, readout simmetrico), calibrati sulle
mediane pubblicate di \texttt{ibm\_torino}, condivisi da tutta la
pipeline a valle.
\end{okbox}

\needspace{4\baselineskip}
\section*{File prodotti}
\addcontentsline{toc}{section}{File prodotti}

\texttt{noise\_model\_dimero.py}, questo documento. Derivazione completa
delle formule in \texttt{teoria\_modello\_rumore.tex}.

\end{document}
